(1) 角函数 \(\cos\theta\) 是空间反演下的奇函数，定 \(j\) 转动波函数宇称为 \((-1)^j\)。因此 \(\langle j,m|\cos\theta|j,m\rangle=0\)，每个无场基态的一阶能移为零。线性极性分子仍有永久电偶极，但自由转动定态的空间平均偶极为零。选择规则 \(\Delta j=\pm1,\Delta m=0\)，可用球谐函数三角积分或 \(3j\) 系数得出。

(2) 对 \(j=0,m=0\)，仅与 \(j=1,m=0\) 耦合，矩阵元 \(\langle1,0|\cos\theta|0,0\rangle=1/\sqrt3\)。二阶非简并微扰给 \(\Delta E_{0,0}^{(2)}=(d\mathcal E)^2(1/3)/(E_0-E_1)=-(d\mathcal E)^2/(6B)\)。由 Hellmann--Feynman 关系 \(p_z=-\partial E/\partial\mathcal E\) 得 \(p_z=d^2\mathcal E/(3B)+O(\mathcal E^3)\)，故静态转动极化率 \(\alpha_{\rm rot}=d^2/(3B)\)。

(3) 对于 \(j=1,m=0\)，向 \(j=0\) 的矩阵元平方为 \(1/3\)，向 \(j=2,m=0\) 的矩阵元平方为 \(4/15\)，对应能量差 \(E_1-E_0=2B,\ E_1-E_2=-4B\)。故 \(\Delta E_{1,0}^{(2)}=(d\mathcal E)^2[(1/3)/(2B)-(4/15)/(4B)]=(d\mathcal E)^2/(10B)\)。对于 \(m=\pm1\)，因 \(j=0\) 无此 \(m\)，仅向 \(j=2,m=\pm1\) 耦合，矩阵元平方 \(1/5\)，所以 \(\Delta E_{1,\pm1}^{(2)}=-(d\mathcal E)^2/(20B)\)。

(4) 因此原本三重简并的 \(j=1\) 能级在二阶分成一条 \(m=0\) 能级和一对仍简并的 \(m=\pm1\) 能级，两者能量差 \(\Delta E_{1,0}-\Delta E_{1,\pm1}=3(d\mathcal E)^2/(20B)\)。当 \(\mathcal E\to0\) 时分裂按电场平方消失。结果需要 \(d\mathcal E\ll B\)，若场强接近转动间隔则应数值对角化跨不同 \(j\) 的耦合，不能继续使用截断的二阶公式。
